\chapter{Adaptations par pilier}
\label{chap:adaptations}

\begin{cle}
\textbf{Idée directrice, et statut de cette annexe.} Le modèle fait passer le Vasicek d'un
facteur unique à un facteur \emph{par pilier}. Les autres composants doivent-ils suivre ? Les
cinq piliers ne sont pas cinq exemplaires du même risque : la queue de sévérité,
l'accumulation par tiers et la détection sont des objets de nature différente. Cette annexe
les adapte pilier par pilier, sous la \emph{même} discipline que le
chapitre~\ref{chap:partielle} : on n'adopte un raffinement que si la donnée le soutient, et
ce qu'elle ne fixe pas, on le borne. Elle est reportée hors du corps parce que son
verdict est constant et ne modifie aucune conclusion : chaque adaptation relève le
niveau, aucune ne renverse le classement.
\end{cle}

\section{La queue de sévérité par pilier}
\label{sec:xi-pilier}

Aujourd'hui une seule queue $\xi$ est appliquée aux cinq piliers. Le
chapitre~\ref{chap:donnees} note pourtant que la queue est hétérogène selon la catégorie
Bâle : $\xi$ varie de $0{,}87$ (\emph{Employment}) à $1{,}37$ (\emph{Execution, Delivery \&
Process Management}), soit une dispersion de $0{,}50$, sur une table recalculée pour ce mémoire.
% SOURCES-SCRIPTS: 67
Faut-il un $\xi_j$ par pilier ? La réponse distingue deux choses.

\begin{attention}
\textbf{L'hétérogénéité est un fait, l'assignation n'en est pas un.} On ne peut pas
\emph{poser} $\xi_j = f(\text{pilier})$ : les catégories Bâle décrivent la nature de la
perte, pas les domaines de contrôle DORA, et la catégorie la plus proche du TIC
(\emph{Business Disruption \& System Failures}) n'a que $111$ observations, sous le seuil
d'estimation de $150$. On applique donc la logique du chapitre~\ref{chap:partielle} : chaque pilier
reçoit une des cinq queues observées, on énumère les $120$ assignations, et l'on borne le
SCR. Pour isoler l'hétérogénéité du niveau, les cinq queues sont centrées sur la valeur de
référence.
\end{attention}

Le résultat est net, et en partie inattendu. L'hétérogénéité de queue n'est pas
neutre en capital : même en gardant la moyenne des queues fixe, donner à chaque pilier sa
propre queue relève le SCR de $+55\,\%$ à $+96\,\%$. C'est un effet de convexité, la
VaR croissant plus vite que linéairement en $\xi$, de sorte qu'une queue lourde sur un pilier
n'est pas compensée par une queue légère ailleurs. Supposer une queue commune
sous-estime donc le capital, une raison de plus de lire le niveau en bande. Le classement,
lui, résiste : le pilier de tête reste dominant dans $83\,\%$ des assignations
(figure~\ref{fig:xi-pilier}), dominant sans être invariant.

\begin{figure}[htbp]\centering
\figover{Z8_xi_par_pilier.png}
\caption{Queue de sévérité par pilier. (a)~Le niveau du SCR borné sur les $120$ assignations
des queues observées aux piliers ; une queue commune (trait) est en bas de la bande.
(b)~Le pilier de tête par contribution à la queue ne dépend pas de l'assignation.}
% SOURCES-SCRIPTS: 37
\label{fig:xi-pilier}
\end{figure}

\section{Le pilier des tiers comme n\oe{}ud d'accumulation}
\label{sec:p4-accumulation}

P4 (prestataires tiers) n'est pas un n\oe{}ud de cascade comme les autres. Une défaillance
d'un prestataire \emph{partagé} dégrade plusieurs piliers simultanément, par
dépendance commune, au lieu de se propager de proche en proche : c'est le mécanisme de MOVEit
($191$ victimes d'un seul choc, chapitre~\ref{chap:donnees}). On lui
% SOURCES-SCRIPTS: 08g
substitue donc un
\emph{choc commun}, où P4 et chaque autre pilier tombent ensemble avec une probabilité
$\varphi_{cs}$, plutôt qu'une propagation dirigée.

\begin{cle}
\textbf{Pourquoi P4 échappe à la limite d'identification de $W$.} Le choc commun est une
\emph{co-occurrence symétrique} : P4 et $P_j$ tombent ensemble, sans direction. Or la
co-occurrence est la partie \emph{identifiée} du modèle, le $S$ de la décomposition
$W=S+A$ (chapitre~\ref{chap:partielle}), celle que la donnée voit (MOVEit, concentration
cloud). L'accumulation de P4 est donc la seule structure inter-piliers que la donnée
soutient, par contraste avec la direction, non identifiée. C'est un point de force, pas une
limite.
\end{cle}

L'intensité $\varphi_{cs}$ au sein d'une entité isolée n'est pas identifiée (le \emph{batch}
observé est un phénomène de portefeuille) : on la borne sur $[0,\gamma]$, où $\gamma=0{,}68$
est la concentration cloud déjà retenue. Le SCR se lit alors en bande, et à
$\varphi_{cs}=\gamma$ une dépendance partagée emporte en moyenne $3{,}7$ piliers d'un coup,
soit $+19\,\%$ de SCR par rapport au traitement en n\oe{}ud de cascade. Surtout, la
simultanéité crée une co-occurrence multiple que la propagation séquentielle ne produit
presque jamais : la probabilité qu'un incident P4 touche au moins trois piliers passe de
$0{,}22$ (cascade) à $0{,}90$ (choc commun), figure~\ref{fig:p4}.

\begin{figure}[htbp]\centering
\figover{Z9_p4_accumulation.png}
\caption{Le pilier des tiers en accumulation. (a)~SCR borné sur l'intensité du choc partagé
$\varphi_{cs}\in[0,\gamma]$, comparé au traitement de P4 en n\oe{}ud de cascade (trait).
(b)~La simultanéité crée une co-occurrence multiple, structure symétrique donc identifiée.
}
% SOURCES-SCRIPTS: 38
\label{fig:p4}
\end{figure}

\subsection{Le seuil de déclenchement, et l'accumulation bornée par la dépendance de queue}
\label{sec:p4-seuil}

Il faut séparer deux questions, celle de la \emph{marge} et celle de la \emph{dépendance}.

\paragraph{La marge : quand un pilier bascule.} Le déclenchement est un franchissement,
\begin{equation}
  \{\text{pilier } j \text{ déclenché}\} \;=\; \{X_j \ge K_j\},
  \qquad K_j = F^{-1}(1-p_j),
  \label{eq:seuil-Kj}
\end{equation}
et pour $P_4$, $p_4=0{,}151$ donne $K_4=1{,}03$. La loi $F$ de la latente n'y est qu'une
\emph{normalisation} : par construction $\Prob(X_j\ge K_j)=p_j$ \emph{quelle que soit} $F$
continue et strictement croissante, de sorte que le choix de $F$ ne transporte aucune information
sur la fréquence de déclenchement. La queue lourde du capital vit dans la sévérité (GPD,
chapitre~\ref{chap:socle}), non dans ce seuil : une loi normale y est donc inoffensive.

\paragraph{La dépendance : combien de piliers ensemble.} L'accumulation, elle, est gouvernée par
la dépendance de queue supérieure de la copule $C$ liant les latentes,
\begin{equation}
  \lambda_U \;=\; \lim_{u\to 1^-}\Prob\bigl(U_2>u \,\big|\, U_1>u\bigr)
             \;=\; \lim_{u\to 1^-}\frac{1-2u+C(u,u)}{1-u} .
  \label{eq:lambda-U}
\end{equation}
Une copule gaussienne vérifie $\lambda_U=0$ quelle que soit sa corrélation, tandis que la copule
de Gumbel de paramètre $\theta>1$ donne
\begin{equation}
  \lambda_U^{\text{Gumbel}} \;=\; 2-2^{1/\theta},
  \label{eq:lambda-gumbel}
\end{equation}
soit $0{,}53$ pour le $\theta=1{,}8$ retenu ailleurs dans le mémoire. Retenir une copule
gaussienne ici sous-estimerait donc structurellement le co-extrême, en contradiction avec ce
choix.

\begin{attention}
\textbf{Borner l'accumulation, à dépendance moyenne égale.} On ne tranche pas la copule : on
lit le co-déclenchement en bande, à $\tau$ de Kendall égal, entre le plancher gaussien
et le haut de Gumbel. La fréquence \emph{marginale} d'un pilier reste $p_4=0{,}15$ pour les deux
(la copule ne change pas combien un pilier bascule seul, mais s'ils basculent \emph{ensemble}) ;
seul le co-extrême diverge : la probabilité que les cinq piliers tombent avec P4 passe
de $0{,}17$ (gaussien) à $0{,}39$ (Gumbel), un facteur $2{,}3$, exactement le scénario MOVEit
que la queue gaussienne manque (figure~\ref{fig:p4-seuil}). Le co-déclenchement d'un autre
pilier tient alors dans $[0{,}48\,;\,0{,}58]$, sous la borne $\gamma=0{,}68$ du choc commun, qui
reçoit ainsi son origine structurelle. Tout est \emph{par pilier} : un nombre de piliers de
l'entité, non une fraction de portefeuille : la vue inter-clients (MOVEit et ses centaines
d'entités victimes) serait une loi de Vasicek en dimension entité, hors du SCR par entité.
\end{attention}

\begin{figure}[htbp]\centering
\figover{Z12_seuil_declenchement_p4.png}
\caption{Le seuil de P4 et l'accumulation bornée. (a)~la marge : franchissement de $K_4$, où la
normale n'est qu'une normalisation. (b)~à dépendance moyenne égale, le co-dépassement s'effondre
vers $0$ pour la gaussienne mais plafonne à $0{,}53$ pour la Gumbel. (c)~l'accumulation par
pilier, la Gumbel concentrant la masse sur les cinq piliers groupés.}
% SOURCES-SCRIPTS: 42
\label{fig:p4-seuil}
\end{figure}

\subsection{Le déclenchement comme temps d'arrêt}
\label{sec:p4-dynamique}

Le seuil $z^*$ est jusqu'ici un niveau statique. Si le facteur tiers est un processus
$(Z_t)_{t\ge0}$ adapté à une filtration $(\mathcal{F}_t)$, le déclenchement devient un
temps de premier passage
\begin{equation}
  \tau \;=\; \inf\{\,t\ge 0 \;:\; Z_t \ge z^*\,\},
  \label{eq:temps-arret}
\end{equation}
qui est un temps d'arrêt, puisque $\{\tau\le t\}=\{\sup_{s\le t}Z_s\ge z^*\}\in
\mathcal{F}_t$. Le \og{}moment\fg{} de déclenchement prend ainsi un sens temporel littéral.

\begin{proposition}[Loi du délai de déclenchement]\label{prop:premier-passage}
Si $(Z_t)$ est un mouvement brownien standard issu de $0$, donc une martingale, alors pour tout
horizon $T>0$
\begin{equation}
  \Prob(\tau\le T) \;=\; 2\,\Phi\!\left(-\frac{z^*}{\sqrt{T}}\right)
  \;=\; 2\,\Prob(Z_T\ge z^*).
  \label{eq:reflexion}
\end{equation}
\end{proposition}

\begin{proof}
C'est le principe de réflexion. Par symétrie du mouvement brownien après $\tau$, on a
$\Prob(\tau\le T,\;Z_T< z^*)=\Prob(\tau\le T,\;Z_T> z^*)=\Prob(Z_T>z^*)$, la dernière égalité
tenant car $\{Z_T>z^*\}\subset\{\tau\le T\}$. En sommant les deux événements disjoints,
$\Prob(\tau\le T)=2\,\Prob(Z_T>z^*)=2\,\Phi(-z^*/\sqrt{T})$.
\end{proof}

La lecture dynamique vaut donc \emph{le double} de la lecture statique $\Prob(Z_T\ge z^*)$ : elle
compte tous les franchissements survenus avant l'horizon, et non la seule position finale. À
$\rho=0{,}50$, la probabilité de déclenchement de P4 en moins d'un an est $0{,}14$ (délai médian
de l'ordre de cinq ans) ; à $\rho=0{,}75$, elle monte à $0{,}23$ (médian d'environ trois ans),
figure~\ref{fig:p4-dynamique}.

\begin{figure}[htbp]\centering
\figover{Z16_p4_temps_arret.png}
\caption{Le déclenchement de P4 comme temps d'arrêt. (a)~le facteur tiers (martingale)
et le premier passage de $z^*$. (b)~la loi du délai de déclenchement (premier passage).
\textbf{(c)}~la lecture dynamique $\Prob(\tau\le 1\text{ an})$ vaut environ le double de la
statique.}
% SOURCES-SCRIPTS: 52
\label{fig:p4-dynamique}
\end{figure}

\begin{cle}
\textbf{Un temps de premier passage, et la propriété utilisée n'est pas celle qu'on nomme.}
$Z_t$ sans dérive est une martingale, mais le principe de réflexion ne s'appuie pas sur cette
propriété : il repose sur la propriété de Markov forte et sur la symétrie des
accroissements. Nommer la martingale ici était décoratif, et il faut se garder du rapprochement
avec le chapitre~\ref{ch:identifiabilite}, où le même mot désigne la forme de Dirichlet d'une
chaîne : le lien entre les deux usages est de vocabulaire et non de mathématique. Ce
qu'ils ont réellement en commun est plus modeste et plus important, ils vivent tous deux sous la
mesure physique. Aucun test de martingalité au sens des générateurs de scénarios
économiques n'est conduit dans ce mémoire, aucun changement de mesure n'y intervient, et le
seuil $z^\star$ ne se déduit d'aucun test de ce genre : il vient de la calibration du facteur
tiers. Le cadre complet, avec les hypothèses effectivement requises, est au
\S\ref{sec:ann-martingalisation}.

Réserve assumée : le brownien pur atteint le niveau presque sûrement mais d'espérance de délai
infinie ; un facteur \emph{moyen-réversif} (Ornstein-Uhlenbeck) donnerait un taux de
déclenchement récurrent fini, plus réaliste pour un facteur systémique. Le brownien est le cas
propre qui porte le calcul explicite.
\end{cle}

\subsection{Le sous-process tiers sur donnée ouverte : ce qu'on peut ancrer avant le registre}
\label{sec:p4-sousprocess}

Descendre P4 au niveau \emph{sous-process} (quel sous-traitant, quelle fonction critique, quelle
substituabilité) suppose le registre d'externalisation DORA, dont la structure est fixée par
l'ITS \citep{eu2024_2956} mais que seules les données de l'entité fournissent. En attendant, la
donnée ouverte tranche déjà deux points opposés (figure~\ref{fig:p4-open}). D'un côté, la
\textbf{taxonomie fine} du tiers y est trop rare : la sous-catégorie \og Vendors \& Suppliers \fg{}
d'OpRisk ne compte que $20$ incidents en finance ($1\,\%$ du TIC), et \texttt{actor.Partner} du
VCDB $56$, qui s'effondrent à zéro deux niveaux plus bas : le même mur de sparsité que partout,
qui renvoie au registre. De l'autre, l'origination et la concentration tierces
sont, elles, ancrées sur donnée européenne officielle, ce qui \emph{valide de l'extérieur} deux
paramètres jusqu'ici posés :
\begin{itemize}\itemsep2pt
\item le premier rapport d'incidents DORA des ESAs \citep{esas2026_incidents} établit que
      $29\,\%$ des $3\,383$ incidents \textsc{tic} majeurs de $2025$ proviennent d'une
      défaillance de tiers : P4 est bien une amorce majeure, ce que le classeur posait
      par jugement ($\text{\texttt{ROOT}}[4]=0{,}90$) ;
\item l'analyse des registres d'externalisation par la BCE \citep{ecb2024_outsourcing} montre que
      plus de $30\,\%$ du budget des grandes banques va à $10$ fournisseurs (et $50\,\%$ du budget
      critique à $30$), et les ESAs ont désigné $19$ fournisseurs \textsc{tic} critiques : cela
      ancre le paramètre d'accumulation $\gamma=0{,}68$, jusque-là un proxy cloud, sur un
      chiffre de concentration officiel.
\end{itemize}

\begin{figure}[htbp]\centering
\figover{Z14_p4_sousprocess_open.png}
\caption{Le sous-process tiers sur donnée ouverte. (a)~la taxonomie fine du tiers est sous le
seuil d'estimation (registre DORA requis). (b)~la concentration, elle, est nette, côté exposition
(BCE) comme côté sinistralité (Gini $0{,}69$). (c)~origination et
concentration tierces ancrées sur donnée UE officielle.}
% SOURCES-SCRIPTS: 44 08g
\label{fig:p4-open}
\end{figure}

\section{La détection de P2 et P3 comme courbe de survie}
\label{sec:detection}

P2 (gestion des incidents) et P3 (tests) ne sont pas des n\oe{}uds de sévérité : ils
gouvernent le temps de détection. Un incident non détecté à temps escalade et
devient matériel. Le modèle résume cela par un multiplicateur du taux de matérialité
($0{,}85$ conforme, $1{,}20$ non conforme). On le remplace par une lecture de survie,
ancrée sur donnée.

Le VCDB donne la distribution du délai de découverte en finance ($n=278$), et elle est
\emph{bimodale} : $47\,\%$ le jour même, $8\,\%$ en semaines, $45\,\%$ en mois ou plus. La
matérialité vient des incidents à découverte lente, d'où un taux de matérialité proportionnel
à la part de découverte lente, ancrée à $\pi_{\text{slow}}=0{,}45$.

\begin{attention}
\textbf{Un point observé, une élasticité bornée.} Le point $\pi_{\text{slow}}=0{,}45$ est
mesuré ; l'élasticité de la détection à la maturité de P2/P3 (de combien un pilier conforme
détecte-t-il plus vite ?) ne l'est pas, et on la borne. Le surcoût du \emph{seul} canal
détection est alors une bande, de $1\,775$ à $8\,052$~M\euro{} selon l'élasticité. Fait
notable : le multiplicateur grossier actuel ($0{,}85/1{,}20$) est retrouvé à une élasticité
$\eta\approx 0{,}36$ : il n'est donc pas arbitraire, c'est un point de la courbe de survie
ancrée sur le VCDB (figure~\ref{fig:detection}).
\end{attention}

\begin{figure}[htbp]\centering
\figover{Z10_detection_survie.png}
\caption{P2/P3, la détection comme courbe de survie. (a)~La distribution bimodale observée
au VCDB. (b)~Le multiplicateur de matérialité comme point de la courbe de survie ; le
multiplicateur grossier actuel correspond à $\eta\approx0{,}36$. (c)~Le surcoût du canal
détection, borné sur l'élasticité.}
% SOURCES-SCRIPTS: 39 67
\label{fig:detection}
\end{figure}

\section{Cas d'usage : des adaptations aux KPI pilotables}
\label{sec:kpi}
% SOURCES-SCRIPTS: 43 68 82 88

Les adaptations qui précèdent se lisent enfin comme un tableau de bord opérationnel, opposable à
un souscripteur ou à un superviseur. Chaque KPI observable, de la fréquence d'entrée au délai de
notification (P2), au délai de confinement (P3) et à la concentration tiers (P4), est rattaché à un
canal du moteur, et l'on mesure le capital en jeu derrière lui : l'écart de capital entre l'état
conforme et l'état non conforme, soit $14\,139$~M\euro{} ($6\,049 \to 20\,188$).

\paragraph{Un canal se lit de trois façons, et elles ne sont pas interchangeables.} On peut partir
de l'état conforme et n'y dégrader que ce canal (lecture \emph{isolée}), partir de l'état non
conforme et n'y remédier que ce canal (marginal de \emph{fermeture}), ou moyenner sur les
vingt-quatre ordres d'entrée (valeur de \emph{Shapley}). Les trois colonnes sont publiées au
tableau~\ref{tab:shapley} et une seule somme à l'écart total : la lecture isolée manque
$5\,001$~M\euro{} et celle de fermeture le dépasse d'autant, le rapport des deux allant de
$1{,}5$ à $3{,}1$ selon le canal. Le sens de cet écart est le point utile pour une direction des
risques : remédier un canal rapporte davantage à une entité défaillante partout qu'à une entité
déjà conforme, puisque ce canal ferme aussi les effets croisés qu'il portait. Seule la lecture de
fermeture correspond donc à une entité réelle en début de remédiation.

\begin{attention}
\textbf{Deux des trois sommes ne mesurent rien, et il ne faut pas les resommer.} La question a été
posée en séance sous la forme d'un soupçon de double comptage, la somme des fermetures paraissant
proche du double de celle des effets isolés. Il n'y a pas de double comptage, et c'est une
identité : en notant $m(S)$ les coefficients de Möbius,
\[
  \sum_i \text{isolé}(i) = \!\!\sum_{|S|=1}\!\! m(S), \qquad
  \sum_i \text{fermeture}(i) = \sum_S |S|\,m(S), \qquad
  \sum_i \varphi_i = \sum_S m(S) = \Delta .
\]
Fermer un canal lui fait porter l'intégralité des croisés auxquels il participe, et un croisé
d'ordre $k$ est donc compté $k$ fois :
$1\times 9\,138 + 2\times 5\,345 + 3\times(-691) + 4\times 346 = 19\,141$, à la précision machine.
Ce qui est en revanche une propriété est l'encadrement $9\,138 \leq 14\,139 \leq 19\,141$, valable
dès que la fonction de capital est sur-modulaire, ce que la cascade est. Shapley est la seule
lecture qui répartisse l'interaction sans la perdre ni la dupliquer, mais c'est une
\emph{convention} d'attribution et non une prévision : deux des quatre canaux sont bornés, de
sorte que leur part n'est pas un budget actionnable. Le défaut de la colonne isolée est celui que
l'analyse de sensibilité rencontre dès que les entrées cessent d'être indépendantes, les indices
de Sobol ne sommant alors plus au total quand la valeur de Shapley y somme par construction
\citep{Owen2014, IoossPrieur2019} ; ces travaux décomposent toutefois une \emph{variance} quand la
fonction caractéristique employée ici est un \emph{écart de capital}, donc l'argument se transpose,
pas le résultat.
\end{attention}

\paragraph{L'état non conforme est le \emph{coin supérieur} du pavé, et c'est une propriété
testée.} Sur les soixante-cinq paires emboîtées des seize configurations, relâcher un canal de plus
ne fait jamais baisser le capital, graine par graine et non seulement en moyenne. Le maximum est
donc atteint au coin haut, et un test de résistance sur le pavé des quatre canaux se réduit à une
seule évaluation. C'est la forme que prend ici le résultat de statique comparative monotone du
préprint cité au chapitre~\ref{ch:positionnement} \citep{ccrn2026}.

\begin{attention}
\textbf{Deux réserves, et la seconde commande la lecture du chiffre de tête.} La borne vaut au sens
du quantile et non trajectoire par trajectoire : à uniformes communs, la perte d'une année donnée
baisse dans $16{,}2\,\%$ des cas quand on relâche la propagation et dans $30{,}7\,\%$ quand on
relâche la détection, pour un motif qui diffère selon le canal, la table des sous-ensembles dans le
premier cas et la transformation de sévérité dans le second. Obtenir la version trajectorielle
demanderait un couplage monotone, donc d'autres tirages, ce que le gel interdit pour un
renforcement d'énoncé sans déplacement de conclusion. Et un coin peut être \emph{infaisable} :
deux des quatre canaux sont des bornes posées et non des valeurs observées, donc rien ne garantit
qu'une entité présente les quatre canaux à leur maximum simultanément. Le coin est un majorant sur
un pavé déclaré, non la description d'une entité.
\end{attention}

% =====================================================================
\subsection{Le test de résistance inversé}
\label{sec:inverse}
% SOURCES-SCRIPTS: 68 84 88 92

La monotonie qui précède autorise la question réciproque, celle qu'un dispositif ORSA pose aussi :
au lieu de fixer un état et de lire un capital, on fixe le \emph{capital} et l'on cherche les états
qui le produisent. L'ensemble des configurations atteignant une cible est alors \emph{croissant},
donc entièrement décrit par ses éléments minimaux, qui ne sont jamais plus de quatre sur seize.

\begin{center}\footnotesize
\renewcommand{\arraystretch}{1.2}
\begin{tabular}{@{}r l l@{}}
\toprule
\textbf{Cible} & \textbf{configurations minimales qui l'atteignent} & \textbf{minimum de canaux} \\
\midrule
$8\,000$ & fréq. ; dét.$+$prop. ; dét.$+$accum. ; prop.$+$accum. & $1$ \\
$10\,000$ & fréq. ; dét.$+$prop.$+$accum. & $1$ \\
$12\,000$ & fréq.$+$dét. ; fréq.$+$prop. ; fréq.$+$accum. & $2$ \\
$15\,000$ & les trois triplets contenant la fréquence & $3$ \\
$18\,000$ & les quatre canaux & $4$ \\
\bottomrule
\end{tabular}
\end{center}

\vspace{4pt}
\noindent \textbf{Deux nombres résument le tableau.} Le canal de fréquence seul atteint
$10\,377$~M\euro, soit $1{,}72$ fois l'état conforme, quand les trois autres canaux réunis
n'atteignent que $11\,413$, soit $1{,}89$ fois, et l'ordre de ces deux bornes tient sur les quatre
graines. Toute cible sous le premier seuil se produit donc par un canal unique ; toute cible
au-dessus du second exige la fréquence. Aucun autre canal seul n'atteint même la cible la plus
basse, la détection plafonnant à $7\,497$~M\euro, la propagation à $7\,682$ et l'accumulation à
$7\,777$ : une exposition à un canal unique est, dans ce modèle, une exposition à la fréquence.

\begin{attention}
\textbf{Le canal d'accumulation n'admet pas d'inversion continue, et son effet isolé est un
\emph{net}.} À l'état conforme, le pilier des tiers est un n\oe{}ud de cascade ordinaire ; à l'état
non conforme il devient un choc commun. Ce sont deux structures de table et non les extrémités d'un
continuum, si bien qu'un test inverse ne peut pas demander à ce canal jusqu'où il doit se dégrader.
La bascule à paramètre nul fait d'ailleurs baisser le capital de $239$~M\euro, résolu sur les
quatre graines, parce qu'elle retire la propagation propre du pilier avant d'ajouter le choc
commun : l'effet isolé publié de $1\,728$~M\euro{} est le net d'un retrait de $239$ et d'un ajout
de $1\,967$, une raison de plus de ne jamais sommer les quatre canaux.
\end{attention}

\noindent Deux réserves, de nature différente. L'inatteignabilité porte sur un pavé \emph{déclaré}
et jamais sur une entité, deux canaux sur quatre étant des bornes posées. Et une configuration
proche d'une cible n'est pas classée, son appartenance dépendant de la graine : une frontière non
résolue est un résultat, elle dit que la cible tombe dans le bruit de simulation.

\paragraph{L'interaction n'est plus un résidu : elle se décompose, et elle vit aux paires.} La
décomposition de Möbius sur le treillis des seize configurations rend compte des $+5\,001$~M\euro{}
terme par terme : quatre effets principaux ($9\,138$~M\euro), six croisés de paires ($5\,345$),
quatre de triplets ($-691$) et un croisé des quatre ($+346$), dont la somme vaut $14\,139$ à la
précision machine. Ce n'est pas une vérification statistique mais une identité : elle atteste que
la table est complète. Les trois paires porteuses passent toutes par la fréquence et font
$86\,\%$ de l'ordre~2, fréquence\,$\times$\,détection $1\,739\pm708$,
fréquence\,$\times$\,accumulation $1\,666\pm354$ et fréquence\,$\times$\,propagation
$1\,182\pm444$~M\euro. Ces $\pm$ sont des étendues de \emph{simulation} à paramètres fixés, et une
loi d'échelle le montre : l'écart-type tombe de $1\,785$ à $339$~M\euro{} lorsque le nombre
d'années simulées passe de $10\,000$ à $160\,000$, soit une pente de $-0{,}60$ contre $-0{,}50$
attendu pour du Monte-Carlo. Cette largeur s'achète en temps de machine ; celle qui ne s'achète pas
est dix fois plus grande, le même croisé allant de $429$ à $7\,998$~M\euro{} aux deux bornes de
l'\textsc{ic}90 de $\hat\xi$. Le classement des deux premières paires n'est pas résolu et n'a pas à
l'être, elles se tiennent à $73$~M\euro{} et leur ordre s'inverse en perte moyenne.

\paragraph{Les états mixtes décomposent, ils ne se reportent pas.} Sur les seize configurations,
quatorze mélangent des canaux conformes et non conformes, et aucune entité ne se trouve dans un tel
état. Ce sont ici des intermédiaires de calcul : ils répartissent un total dont les deux extrémités
sont cohérentes, et la décomposition les fait disparaître par construction. Un chiffrage qui
\emph{reporterait} un état mixte présenterait comme atteignable une situation qui ne l'est pas.

\paragraph{Une seule paire est négative, et elle borne ce qu'on peut attribuer à la contagion.}
Propagation\,$\times$\,accumulation vaut $-330\pm242$~M\euro, négative sur les seize graines et sur
les deux métriques. Les deux canaux sont substituts et non compléments : la cascade et le choc
commun sont deux façons de faire tomber plusieurs piliers sur un même sinistre, et le premier
actif, il reste moins à gagner au second. Le capital est donc super-additif en euros mais
sous-additif entre ses deux canaux de co-occurrence, ce qui borne l'amplitude imputable à la
contagion prise en général.

\paragraph{D'où vient l'interaction : les canaux se composent de façon quasi multiplicative.} Une
décomposition \emph{additive} d'effets qui se \emph{multiplient} produit une interaction positive
par simple arithmétique, $(1+a)(1+b)-1$ dépassant $a+b$. La décomposition refaite sur le logarithme
du capital le montre. Le produit des quatre facteurs isolés vaut $3{,}47$ quand le facteur publié
vaut $3{,}34$, soit $0{,}961$ fois ce produit ; un modèle purement multiplicatif produirait $5\,813$~M\euro{}
d'interaction, et les $5\,001$~M\euro{} publiés en valent $86\,\%$. Aucune des seize configurations
ne s'écarte de plus de $8{,}1\,\%$ du produit de ses facteurs isolés. L'interaction en euros est
donc entièrement expliquée par la multiplicativité, et le résidu propre est une légère
\emph{sous}-multiplicativité, portée par la paire propagation\,$\times$\,accumulation. Deux canaux
suivent même une forme fermée : par l'approximation de la perte unique
\citep{bockerkluppelberg2005}, le quantile croît comme
$(\lambda\,p_u)^{\xi}$, ce qui prédit un facteur $(53{,}62/21{,}56)^{0{,}5954}=1{,}720$ pour la
fréquence contre $1{,}715$ simulé, et $1{,}228$ pour la détection contre $1{,}239$. Les conclusions
de gestion n'en changent pas, remédier un canal rapportant toujours davantage à une entité
défaillante partout ; elles changent de statut, puisqu'elles se démontrent au lieu de se constater,
et l'interaction ne se lit pas comme un effet propre de la contagion.
% SOURCES-SCRIPTS: 100

\begin{figure}[htbp]\centering
\figover{S24_interaction_canaux.png}
\caption{La table complète des quatre canaux. (a)~la réconciliation par ordre : les quinze termes
de Möbius somment à l'écart, exactement. (b)~les trois lectures d'un même canal, isolée, Shapley
et fermeture. (c)~les six croisés de paires avec leur écart-type de simulation ; une seule est
négative.}
\label{fig:interaction-canaux}
\end{figure}

\begin{cle}
\textbf{Classer par capital \emph{et} par statut.} Les KPI à fort capital et calibrables
(fréquence, détection) sont les leviers de remédiation à activer en priorité ; ceux à fort capital
mais bornés (propagation, accumulation) mesurent surtout une incertitude, à résorber par la donnée,
c'est-à-dire par le registre DORA, plutôt que par une action directe (figure~\ref{fig:kpi}).
\end{cle}

\begin{figure}[htbp]\centering
\figover{Z13_kpi_pilotables.png}
\caption{Des KPI descriptifs aux KPI pilotables. (a)~l'écart de capital DORA attribué par KPI,
chaque canal isolé (bleu : calibrable ; gris : borné), avec l'interaction qui referme la
cascade. (b)~le tableau de bord de remédiation, les KPI classés par capital libéré.}
\label{fig:kpi}
\end{figure}

% =====================================================================
\subsection{Le retour sur investissement de la conformité}
\label{sec:roi}
% SOURCES-SCRIPTS: 50 58 83

Le tableau de bord se lit enfin en euros, et deux mots doivent y être définis avant tout calcul.
Le \emph{portage} d'abord : détenir du capital n'est pas gratuit, et sous Solvabilité~II son coût
de détention se lit dans la marge de risque, un pourcentage par an du capital immobilisé. Au taux
de $6\,\%$ du règlement délégué, les $14\,139$~M\euro{} que libère la conformité totale valent donc
$848$~M\euro{} par an, et non une fois, la fréquence en portant à elle seule près de la moitié.

\paragraph{Le plancher porte sur le bénéfice, non sur la durée.} Le bénéfice a deux composantes et
le calcul ci-dessus n'en monétise qu'une : le portage économisé vaut $848$~M\euro/an quand la
sinistralité évitée en vaut $3\,247$, la charge annuelle moyenne passant de $3\,869$ à
$622$~M\euro. La composante laissée de côté vaut $3{,}8$ fois celle qui est retenue, donc le
portage seul minore le bénéfice. Il en résulte que le retour de $6$ à $35$ ans, calculé face à un
coût de remédiation de $25$ à $150$~M\euro{} par grande entité, soit $5$ à $30$~Md\euro{} pour deux
cents entités, est un \emph{majorant} : le compte réel est plus court, jamais plus long. Les deux
composantes réunies donnent $4\,095$~M\euro/an, soit un retour de $1$ à $7$ ans, mais cette
addition mêle un flux de compte de résultat à un coût du capital : c'est une convention de retour
sur investissement, non une sortie du modèle, et le chiffre de référence reste le portage seul avec
le sens de sa borne.

\paragraph{Ce taux n'est pas un critère de décision d'investissement.} Il rémunère l'immobilisation
de fonds propres sous Solvabilité~II ; il n'est ni le coût moyen pondéré du capital de l'entité, ni
son taux d'exigence sur un projet. Une durée de retour exprimée dans cette unité est une grandeur
de comparaison prudentielle, non un test d'opportunité, et la confusion entre les deux est
fréquente en pratique.

\paragraph{Une réserve sur le taux, déclarée plutôt que corrigée.} Le taux de $6\,\%$ est celui du
règlement délégué~2015/35 ; la directive~(UE)~2025/2 le ramène à $4{,}75\,\%$, taux qu'emploient
les calculs d'échelle d'entité. Le portage monétisé tombe alors à $672$~M\euro/an, soit une baisse
de $21\,\%$, et le majorant du retour passe de $35$ à $45$ ans. Aucune conclusion ne bascule, la
durée demeurant un majorant et la composante dominante du bénéfice n'étant pas monétisée ; la
chaîne publiée reste au taux historique, la calibration étant gelée.

\begin{figure}[htbp]\centering
\figover{J6_roi_conformite.png}
\caption{Le ROI de la conformité. (a)~le capital libéré par levier, classé, avec sa valeur
annuelle à $6\,\%$ (bleu : calibrable, gris : borné). (b)~coût de remédiation \emph{one-off}
contre économie \emph{récurrente} de portage du capital ; le retour sur le portage seul est de
$6$ à $35$ ans, et c'est un majorant, la perte évitée n'étant pas monétisée ici.}
\label{fig:roi}
\end{figure}

\begin{cle}
\textbf{La conformité comme levier de capital.} Le bénéfice, le capital libéré, est issu du modèle ;
le coût est externe et incertain, affiché comme tel. La conclusion tient : la conformité n'est pas
un pur coût mais un levier de capital, à prioriser par capital libéré, les leviers calibrables
d'abord.
\end{cle}

% =====================================================================
\section{Ce qui reste en perspective : P1 et P5}
\label{sec:perspectives-pilier}

Deux piliers appelleraient un traitement propre mais \emph{aucune donnée ne le soutient}, et
il serait contraire à toute la démarche de les modéliser par des paramètres non
identifiables. Ils restent donc en perspective.

\begin{itemize}\itemsep2pt
\item \textbf{P1 (gouvernance)} ressemble moins à un n\oe{}ud qu'à un \emph{modificateur
      systémique} : une gouvernance défaillante charge la sensibilité $\rho$ de tous les
      piliers à la fois. Aucune donnée ne relie la qualité de gouvernance à une hausse de
      hasard : à borner, pas à estimer.
\item \textbf{P5 (partage d'informations)} agit plutôt comme un \emph{réducteur de
      fréquence}, l'alerte précoce diminuant les entrées. Là encore, rien ne le mesure.
\end{itemize}

\begin{cle}
\textbf{Le fil de ce chapitre.} Chaque pilier, correctement modélisé, révèle une structure
\emph{localement identifiée} (hétérogénéité de queue, accumulation symétrique, courbe de
détection) que la donnée soutient sur un point et que l'on borne pour le reste. Toutes ces
adaptations relèvent le niveau du capital, et aucune ne renverse le
classement des piliers. C'est, décliné composant par composant, la thèse du mémoire : le
niveau est illustratif, l'ordre et les écarts sont ce qui résiste, et la limite
d'identification se solde partout en une même spécification de reporting.
\end{cle}
